\documentclass{article}
\input{../../../lib.sty}
\usepackage{bbm}

\title{\texttt{fn make\_expr\_count}}
\author{\MichaelShoemateAuthor}

\begin{document}
\maketitle
\contrib

This document proves that the implementation of \rustdoc{transformations/make_stable_expr/expr_count/fn}{make\_expr\_count} in \asOfCommit{mod.rs}{f5bb719}
satisfies its proof definition.

\section{Hoare Triple}
\subsection*{Precondition}
\subsubsection*{Compiler-verified}
\begin{itemize}
    \item Argument \texttt{input\_domain} of type \texttt{WildExprDomain}
    \item Argument \texttt{input\_metric} of type \texttt{L01InfDistance<MI>}
    \item Argument \texttt{expr} of type \texttt{Expr}
    \item Generic \texttt{MI} must implement \texttt{UnboundedMetric}
    \item Const generic \texttt{P} must be of type \texttt{usize}
    \item \texttt{(ExprDomain, L01InfDistance<MI>)} implements \texttt{MetricSpace}
    \item \texttt{(ExprDomain, LpDistance<P, f64>)} implements \texttt{MetricSpace}
    \item \texttt{Expr} implements \texttt{StableExpr<L01InfDistance<MI>, L01InfDistance<MI>>}
\end{itemize}

\subsubsection*{Caller-verified}
None

\subsection*{Pseudocode}

\lstinputlisting[language=Python,firstline=2,escapechar=`]{./pseudocode/make_expr_count.py}

\subsection*{Postcondition}
\begin{theorem}
\validTransformation{(input\_domain, input\_metric, expr, MI, P)}{make\_expr\_count}
\end{theorem}
Starting from line \ref{match-strategy},
\texttt{expr} is analyzed to determine the type of counting query \texttt{strategy} and input to the counting query \texttt{input}.

All preconditions for \texttt{make\_stable} on line \ref{t-prior} are compiler-verified,
therefore by the postcondition \texttt{t\_prior} is a valid transformation.

To prove that the output is a valid transformation,
we must first prove that the transformation on line \ref{count-trans} is valid.

\begin{proof}
    \textbf{Appropriate Output Domain}
    Since the count transformation is not row-by-row,
    line \ref{alignment} disallows this count transformation from being constructed in a row-by-row context,
    satisfying the requirements of ExprDomain.

    By the definition of each of the allowed counting expressions in \ref{match-strategy},
    the resulting output will always contain one series whose name matches the active series' name,
    and the data type of elements in the series is u32,
    which is consistent with \ref{output-series}.
    Each group has one non-null count of type \texttt{u32}; absent counts are filled with zero.
    The output aggregation context on line \ref{output-domain} retains the grouping keys, maximum number of groups and invariant, and sets the maximum group length to one.
    This domain is then used on line \ref{output-domain-usage} to construct the count transformation.
\end{proof}

\begin{proof}
    \textbf{Stability Guarantee}

    To determine if \texttt{input} is row-by-row,
    line \ref{is-row-by-row} checks if \texttt{input} can be parsed into a row-by-row transformation.

    Line \ref{will-count-all} uses domain descriptors to determine if the counting query will count all rows.
    This is the case when the input is row-by-row and the strategy is \texttt{Len},
    or when the input is row-by-row and the strategy is \texttt{Count} on non-nullable data.

    Only in these cases is the input margin's invariant passed to the stability map, on line \ref{invariant}.
    In particular, invariant group lengths imply invariant counts only when all rows are counted.

    The function defined on line \ref{function} applies the counting query to the input expression with the given query strategy.
    The stability of this function is governed by the stability map returned by \rustdoc{transformations/make_stable_expr/expr_count/fn}{counting\_query\_stability\_map},
    because adding or removing one row changes each of the four counts by at most one, including the number of distinct values.
    A missing group has count zero, so the same bound applies to appearing and disappearing groups.
    These facts satisfy the caller-verified preconditions of the stability map.
\end{proof}

Since it has been shown that both \texttt{t\_prior} and the count transformation are valid transformations,
then the preconditions for \texttt{make\_chain\_tt} are met (invoked via the right-shift operator shorthand).

\end{document}
